\documentclass[10pt,a4paper]{amsart}
\usepackage{amsmath,amssymb,amsthm,graphicx}
% Retained preamble for the Stacks Project chapter snapshot in this directory.
% The chapter's own preamble.tex is a copy of the shared Stacks Project preamble; the block
% below is sliced MECHANICALLY, line for line, from it. Nothing was added, removed or
% redefined; the wrapper class, the xr-hyper/hyperref wiring and the externaldocument list
% are omitted because the assembler supplies the document class.
%
% lines taken from preamble.tex: 3, 5, 9, 11, 12, 156, 157, 158, ... (51 lines total)
%
% package coverage restored from the chapter's own preamble (2026-10-13): xy, xy-2cell

\usepackage[all]{xy}
\xyoption{2cell}
\UseAllTwocells

\newenvironment{reference}{\comment}{\endcomment}
\newenvironment{slogan}{\comment}{\endcomment}
\newenvironment{history}{\comment}{\endcomment}
\theoremstyle{plain}
\newtheorem{theorem}[subsection]{Theorem}
\newtheorem{proposition}[subsection]{Proposition}
\newtheorem{lemma}[subsection]{Lemma}
\theoremstyle{definition}
\newtheorem{definition}[subsection]{Definition}
\newtheorem{example}[subsection]{Example}
\newtheorem{exercise}[subsection]{Exercise}
\newtheorem{situation}[subsection]{Situation}
\theoremstyle{remark}
\newtheorem{remark}[subsection]{Remark}
\newtheorem{remarks}[subsection]{Remarks}
\numberwithin{equation}{subsection}
\def\lim{\mathop{\mathrm{lim}}\nolimits}
\def\colim{\mathop{\mathrm{colim}}\nolimits}
\def\Spec{\mathop{\mathrm{Spec}}}
\def\Hom{\mathop{\mathrm{Hom}}\nolimits}
\def\Ext{\mathop{\mathrm{Ext}}\nolimits}
\def\SheafHom{\mathop{\mathcal{H}\!\mathit{om}}\nolimits}
\def\SheafExt{\mathop{\mathcal{E}\!\mathit{xt}}\nolimits}
\def\Sch{\mathit{Sch}}
\def\Mor{\mathop{\mathrm{Mor}}\nolimits}
\def\Ob{\mathop{\mathrm{Ob}}\nolimits}
\def\Sh{\mathop{\mathit{Sh}}\nolimits}
\def\NL{\mathop{N\!L}\nolimits}
\def\CH{\mathop{\mathrm{CH}}\nolimits}
\def\proetale{{pro\text{-}\acute{e}tale}}
\def\etale{{\acute{e}tale}}
\def\QCoh{\mathit{QCoh}}
\def\Ker{\mathop{\mathrm{Ker}}}
\def\Im{\mathop{\mathrm{Im}}}
\def\Coker{\mathop{\mathrm{Coker}}}
\def\Coim{\mathop{\mathrm{Coim}}}
\DeclareMathSymbol{\boxtimes}{\mathbin}{AMSa}{"02}
\def\QCohstack{\mathcal{QC}\!\mathit{oh}}
\def\Cohstack{\mathcal{C}\!\mathit{oh}}
\def\Spacesstack{\mathcal{S}\!\mathit{paces}}
\def\Quotfunctor{\mathrm{Quot}}
\def\Hilbfunctor{\mathrm{Hilb}}
\def\Curvesstack{\mathcal{C}\!\mathit{urves}}
\def\Polarizedstack{\mathcal{P}\!\mathit{olarized}}
\def\Complexesstack{\mathcal{C}\!\mathit{omplexes}}
\def\Pic{\mathop{\mathrm{Pic}}\nolimits}
\def\Picardstack{\mathcal{P}\!\mathit{ic}}
\def\Picardfunctor{\mathrm{Pic}}
\def\Deformationcategory{\mathcal{D}\!\mathit{ef}}
\usepackage{cleveref}
\crefname{originallemma}{original lemma}{original lemmas}
\Crefname{originallemma}{Original lemma}{Original lemmas}
\makeatletter
\newlabel{lemma-frobenius-homotopy}{{0G5Y}{source}{}{sourcepointer-1}{}}
\newlabel{lemma-frobenius-homotopy@cref}{{[originallemma][0][]0G5Y}{[1][0][]source}{}{}{}}
\makeatother
\crefname{originalsection}{original section}{original sections}
\Crefname{originalsection}{Original section}{Original sections}
\makeatletter
\newlabel{section-functoriality}{{08QL}{source}{}{sourcepointer-2}{}}
\newlabel{section-functoriality@cref}{{[originalsection][0][]08QL}{[1][0][]source}{}{}{}}
\makeatother
\title{The Frobenius-induced endomorphism of the cotangent complex in characteristic p is null homotopic}
\author{The Stacks Project Authors}
\date{stacks-project:cotangent.tex:local-snapshot-20261012; GNU FDL 1.2 or later}
\begin{document}
\maketitle
\setcounter{subsection}{1}
\noindent\textit{Editorial note.} All delivered mathematical text is the author's own, copied from the pinned Stacks Project chapter file cotangent.tex; no statement or proof step was rewritten, reordered or removed.\par\medskip
\noindent\textit{Editorial note.} The statement names the characteristic hypothesis and refers to the functoriality section; the proof is a short homotopy computation.\par\medskip
\noindent\textit{Editorial note.} The two references of the counted spans are recorded as non-load-bearing navigation pointers; the argument that consumes them stays inside the counted proof span.\par\medskip
\noindent\textit{Editorial note.} Preamble: the retained preamble of this chapter (sources/stacks-cotangent/retained-preamble.tex, lines 1-63), containing the Xy-pic setup, the comment environments, the theorem declarations and the macros including \textbackslash{}NL.\par\medskip
\noindent\textit{Allowed standard prerequisites (editorial).}\begin{itemize}
\item The Frobenius homotopy on the standard resolution. Lemma lemma-frobenius-homotopy, used at proof lines 2026-2029. Replaces the given self-map by the termwise Frobenius map up to homotopy.
\item The functoriality of the cotangent complex. Section section-functoriality, used at proof lines 2020-2022. Defines the map whose vanishing up to homotopy is asserted.
\item The vanishing of the derivative of a p-th power. used at proof lines 2030-2032. Shows that the termwise Frobenius acts as zero on the differentials.
\end{itemize}\medskip
\ifdefined\hypertarget\hypertarget{sourcepointer-1}{}\fi
\noindent\textit{Source reference.} Original lemma 0G5Y: focused-seeds/AGP/sources/stacks-snapshot-20261011/tags.txt, source lines 20178--20178. Navigation pointer into the Stacks Project; the standard fact it supplies is applied inside the counted proof span and the argument resting on it stays there.\par\medskip
\ifdefined\hypertarget\hypertarget{sourcepointer-2}{}\fi
\noindent\textit{Source reference.} Original section 08QL: focused-seeds/AGP/sources/stacks-snapshot-20261011/tags.txt, source lines 11064--11064. Navigation pointer into the Stacks Project; the standard fact it supplies is applied inside the counted proof span and the argument resting on it stays there.\par\medskip
\begin{lemma}
\label{lemma-frobenius-acts-as-zero}
Let $p$ be a prime number. Let $A \to B$ be a ring homomorphism
and assume that $p = 0$ in $A$. The map $L_{B/A} \to L_{B/A}$
of Section \ref{section-functoriality} induced by the
Frobenius maps $F_A$ and $F_B$ is homotopic to zero.
\end{lemma}

\begin{proof}
Let $P_\bullet$ be the standard resolution of $B$ over $A$.
By Lemma \ref{lemma-frobenius-homotopy} the map $P_\bullet \to P_\bullet$
induced by $F_A$ and $F_B$ is homotopic to the map
$F_{P_\bullet} : P_\bullet \to P_\bullet$ given by
Frobenius on each term. Hence we obtain what we want as clearly
$F_{P_\bullet}$ induces the zero map $\Omega_{P_n/A} \to \Omega_{P_n/A}$
(since the derivative of a $p$th power is zero).
\end{proof}

\end{document}
